\section*{الفصل \ref{ch:g9:how-species-change} --- كيف تتغير الأنواع}

\begin{solution}{exo:g9:how-species-change:1}
تنوع في \omterm{def:g9:heredity-and-traits:traits}{صفات} موروثة؛ وفرط إنتاج تلقاه كوابح \omterm{def:g6:exploring-our-environment:components}{البيئة}؛ وجمع
غير متساو --- فالملائمون يفلتون أكثر قليلًا ويتركون ورثة
أكثر. والنتيجة: تنمو حصة \omterm{def:g9:genes-and-dna:gene}{الأليلات} الملائمة جيلًا بعد جيل
--- وهو \omterm{prop:g9:how-species-change:selection}{الانتقاء الطبيعي}، ومع الزمن، التطور.
\end{solution}

\begin{solution}{exo:g9:how-species-change:2}
التنوع: صورتان موروثتان، شاحبة وداكنة. والكابح: افتراس
\omterm{prop:g3:sorting-living-things:groups}{الطيور} المبصرة على اللحاء. والجمع غير المتساوي: على
الجذوع المسخّمة ترى الشاحبة فتؤكل --- فترتفع حصة
الداكنة؛ وعلى اللحاء الشاحب النظيف ترى الداكنة ---
فتستعيد الشاحبة حصتها. محرك واحد، وأي اتجاه.
\end{solution}

\begin{solution}{exo:g9:how-species-change:3}
\omterm{def:g8:reproduction-and-environment:population}{الجمهرة}: فلا تسود عثة واحدة. والانتقاء يزيح
\emph{تواترات} الصور الموروثة عبر الأجيال --- فالأفراد
يوزعون وينفقون؛ والذي ينزاح هو تركيب المجموعة.
\end{solution}

\begin{solution}{exo:g9:how-species-change:4}
الجهد لا يورث --- فالعنق الممطوط \omterm{def:g1:my-body:parts}{ذراع} الحداد. والآلية:
الأعناق تنوعت وراثيًا؛ فتغذت الصور الأطول عنقًا أعلى،
وأفلتت من كوابح العجاف أكثر، وتركت ورثة أكثر يحملون
\omterm{def:g9:genes-and-dna:gene}{أليلات} العنق الأطول --- فطالت أعناق \omterm{def:g8:reproduction-and-environment:population}{الجمهرة}، ولم يطل عنق
زرافة واحدة.
\end{solution}

\begin{solution}{exo:g9:how-species-change:5}
المحرك ينقر بالأجيال، وأجيال \omterm{def:g6:producing-our-food:fermentation}{الميكروبات} تدور في ساعات:
فالأيام تسع آلاف النقرات. أضف كابحًا وحشيًا (المضاد
الحيوي) فترث الطافرات المقاومة النادرة الجناح في غضون
الأسبوع.
\end{solution}

\begin{solution}{exo:g9:how-species-change:6}
كلاهما كان الكابح غير المتساوي على تنوع موروث --- يقرر
من يتناسل. والفرق: المربّية اختارت عن قصد نحو غاية؛
\omterm{prop:g3:sorting-living-things:groups}{والطيور} لم تختر شيئًا --- بل انتقت عواقب اصطيادها.
\end{solution}

\begin{solution}{exo:g9:how-species-change:7}
الحقيقة الأولى: فهو لا يقارن إلا الصور \emph{الحاضرة}.
والانتقاء غربال على المجموعة القائمة --- لا يستطيع أن
يستحضر التصميم الأفضل غير المطفَّر وغير المخلوط، بل
يرقّي الأفضل الموجود؛ ومن هنا مخططات \omterm{def:g1:animals-around-us:parts}{أطراف} أعيد
توظيفها، لا هندسة جديدة.
\end{solution}

\begin{solution}{exo:g9:how-species-change:8}
التنوع: عمق المنقار، وهو موروث (تبينه أنساب الباحثين).
والكابح: \omterm{def:g2:from-seed-to-plant:seed}{بذور} الجفاف القاسية --- فالمناقير الصغيرة لا
تكسر ما بقي من الغذاء. والمحاسبة: \omterm{prop:g3:sorting-living-things:groups}{الطيور} العميقة المنقار
تتغذى وتبقى وتتناسل؛ فتصير فراخ الموسم التالي أعمق في
المتوسط --- إذ كسبت \omterm{def:g9:genes-and-dna:gene}{أليلات} العمق حصة. ولا انشطار: \omterm{def:g8:reproduction-and-environment:population}{جمهرة}
واحدة، ونقرة واحدة، مقيسة.
\end{solution}

\begin{solution}{exo:g9:how-species-change:9}
الفصل: البحر --- فالغرقى يتناسلون وحدهم. والتباعد: ظروف
كل جزيرة تجري المحرك على طريقتها، فتراكم أليلاتها هي،
جيلًا بعد جيل. \omterm{def:g3:sorting-living-things:criterion}{والمحك}: طيور أعيد جمعها تخفق في التناسل
--- فبتعريف \omterm{def:g5:biodiversity:species}{النوع}، اكتمل الفرق.
\end{solution}

\begin{solution}{exo:g9:how-species-change:10}
كلاهما يطبق الكابح مهلهلًا. فالجرعة غير المتمَّمة تقتل
القابلة للتأثر وتطلق الناجية --- وقد أغنيت بالمقاومة ---
في منتصف الانتقاء؛ والجرعات المنثورة في المزارع تجري
الغربال الجزئي نفسه على \omterm{def:g8:reproduction-and-environment:population}{جمهرات} ميكروبية هائلة يوميًا.
والكابح الذي يستبعد ولا يتم برنامج تربية لما يعفيه.
\end{solution}

\begin{solution}{exo:g9:how-species-change:11}
أثبت فصل الكوابح أن أكثر فائض كل جيل يجمع --- افتراسًا،
وجوعًا، وظروفًا؛ وأثبت فصل الفرادة أن المجموعين يختلفون
وراثيًا، مجموعةً مجموعة. والانتقاء ليس شيئًا مضافًا: إنه
ملاحظة أن تلك الكوابح، إذ تقع على ذلك التنوع، لا يمكن أن
تقع بالسواء --- الفصلان مضروب أحدهما في الآخر.
\end{solution}

\begin{solution}{exo:g9:how-species-change:12}
الانتقاء ينفق التنوع؛ والخلط \omterm{prop:g9:genes-and-dna:explains}{والطفرة} يعيدان تموينه.
\omterm{def:g8:reproduction-and-environment:population}{وجمهرة} المستنسخات تعرض على الغربال مجموعة واحدة: فإذا
انقلبت الظروف لم يجد الانتقاء ما يرقّيه --- فالكرم يقوم
أو يسقط \omterm{def:g1:plants-around-us:plant}{كنبتة} واحدة. ومن غير رشح التهجئات الجديدة وتوزيع
التركيبات الجديدة، يتوقف المحرك: فالتنوع ليس ناتجًا
عرضيًا للآلية، بل هو وقودها.
\end{solution}

\begin{solution}{pb:g9:how-species-change:1}
\textbf{1.} فصل \omterm{def:g6:how-plants-spread:dispersal}{الانتثار}: فلا يبلغ الجزيرة المحيطية إلا
خدمات المدى البعيد --- الأجنحة، والطفو، وما تحمله
العواصف؛ فقائمة الركاب هي بيان تلك الخدمات.

\textbf{2.} أن عصافير الأرخبيل منحدرة من غرقى البر ---
\omterm{def:g9:heredity-and-traits:traits}{فالصفات} المشتركة في التفاصيل تعني انحدارًا مشتركًا،
وأقرب أقربائها يسمي \omterm{def:g8:reproduction-and-environment:population}{الجمهرة} المصدر.

\textbf{3.} فصل \omterm{def:g4:adapted-to-their-world:adaptation}{التكيف} يقرأ الملاءمة --- فكل منقار يجيب
مسألة غذاء جزيرته. وهذا الفصل يضيف \emph{الحدوث}: تنوع
الغرقى، وكوابح الجزيرة، والجمع غير المتساوي، وأجيال من
المحاسبة --- أي صناعة الملاءمة، حيث لم نكن نرى من قبل
إلا الملاءمة (الجواب الذي وعد به
\cref{rem:g4:adapted-to-their-world:how}).

\textbf{4.} الفصل: \omterm{def:g8:reproduction-and-environment:population}{فالجمهرات} المحوطة بالبحر تتناسل
وحدها منذ البداية، وكل جزيرة غرفة محرك مغلقة ---
فالجغرافيا توفر مكوّن الانتواع الأول بالجملة.

\textbf{5.} يصل الكابح (\omterm{def:g2:from-seed-to-plant:seed}{بذور} قاسية فقط)؛ فيجمع التنوع
الموروث في عمق المنقار على غير سواء (فتجوع المناقير
الصغيرة)؛ وتناسل الناجين يكتب الانزياح في متوسط الجيل
التالي --- نقرة واحدة من المحرك، موثقة بالفرجار.

\textbf{6.} أن ``اتجاهه'' ليس أبدًا إلا الظروف القائمة:
فالانتقاء يتتبع الكابح، والكابح المعكوس يعكس الانزياح
--- لا مقصد، بل عواقب، موسمًا بعد موسم.

\textbf{7.} \Cref{prop:g9:how-species-change:speciation}
في منتصف جريانها: تباعد بلغ من التقدم حدًا يتعثر عنده
التناسل --- فرق ينغلق أمام \omterm{def:g1:the-five-senses:sense}{أعين} الدفاتر.

\textbf{8.} نقرات الانتواع أجيال بالآلاف --- ولا ترى
كاملة إلا في حياة سريعة الدوران (\omterm{def:g6:producing-our-food:fermentation}{ميكروبات} الجناح). أما
في العصافير فيقوم المنهج شاهدًا بدلًا من ذلك: النمط،
\omterm{prop:g6:classification-kinship:kinship}{والقرابة}، والنقرات الحية، والفروق نصف المنغلقة تتقارب
على عملية لا يسعها عمر واحد.

\textbf{9.} (1) التنوع: مؤسسون غرقى، مخلوطون
ومطفَّرون. (2) الكوابح: أغذية كل جزيرة وأخطارها. (3)
المحاسبة: \omterm{def:g9:genes-and-dna:gene}{أليلات} كل جزيرة الملائمة تكسب، جزيرة جزيرة.
(4) الفصل: البحر بينها، يمسك كل دفتر على حدة. (5) مقياس
الزمن: آلاف الأجيال --- وقد قامت الفروق تامة.

\textbf{10.} لا شيء إلا ما حمله مؤسسوها وما طفر عندها
منذئذ: فمجموعة كل جزيرة عينة غريق، وما كان الانتقاء
ليرقّي منها إلا أفضل منقار \emph{متاح} --- فالجزر ذوات
المجموعات الأغنى أو الطفرات الأسعد بلغت أجوبة أقرب.

\textbf{11.} زمن الجيل: ساعات في الجناح، وسنوات في
الجزر --- الحقائق الثلاث نفسها، والمحاسبة نفسها، ونقرات
تفصل بينها آلاف الأضعاف.

\textbf{12.} كلاهما كابح يقرر من يتناسل: فيد المربّي
المتعمدة \omterm{def:g2:from-seed-to-plant:seed}{وبذور} الجفاف القاسية العمياء تجريان المحاسبة
نفسها على تنوع موروث --- ولهذا كانت التغيرات التي صنعت
في أبراج الحمام في قرون تصدّق ما تستطيع كوابح الطبيعة أن
تكتبه في ملايين السنين.
\end{solution}
